\documentclass[12pt]{article}
\usepackage[utf8]{inputenc}
\usepackage[T1]{fontenc}
\usepackage{hyperref}
\title{Fala Language Resources for K-12 Curriculum Development}
\author{Compiled by Claude Sonnet 5.5 (claude-sonnet-5-5) for Andr\'as Kornai}
\date{}
\begin{document}
\maketitle

\section{Introduction}
Fala (A Fala, ISO 639-3 \texttt{fax}) is spoken in three villages of the Valle de Xàlima in northwestern C\'aceres, Spain: Eljas (\emph{lagarteiru}), San Mart\'in de Trevejo (\emph{ma\~negu}) and Valverde del Fresno (\emph{valverde\~nu}). Search results indicate that the three varieties are mutually intelligible and that the language is absent from official curricula, although individual school initiatives have brought it into classrooms. Only a few resources were found; this report does not pad beyond them.

\section{Primary Educational Resources}
No dedicated primer, textbook or school curriculum in Fala was found in the searches. The bilingual dictionary below is the closest classroom-usable tool.
\begin{itemize}
\item Diccionariu de A Fala (Fala--Spanish and Spanish--Fala, lagarteiru, ma\~negu and valverde\~nu varieties; second web edition, November 2022): \url{https://afala.fp.tul.cz/images/download/DICCIONARIU_A_FALA_version02_web_Nov2022.pdf}
\end{itemize}

\section{Linguistic and Reference Materials}
\begin{itemize}
\item A Fala project site (English pages), hosted by a Czech university project: \url{https://afala.fp.tul.cz/en/a-fala}
\item Wikipedia overview of the Fala language: \url{https://en.wikipedia.org/wiki/Fala_language}
\item Academic article ``An\'alisis del l\'exico diferencial en el Val de X\'alima (C\'aceres)'': \url{https://portalcientifico.uvigo.gal/documentos/5edec7c22999526e5fc99294}
\end{itemize}

\section{Religious and Community Resources}
No religious texts or community organisation pages for Fala were found in the searches.

\section{Digital and Audio Resources}
No audio or video collections were found in the searches. The dictionary PDF above is the only digital resource located.

\section{Curriculum Development Context}
Each of the three villages has an elementary school and Valverde del Fresno also has a middle school, but Fala is not part of the official curriculum. A curriculum would need to start from the dictionary and the project site, and would require new primers and readers, with the three local varieties kept distinct.
\end{document}